% from ega0-1-fr.tex
\oldpage[0\textsubscript{I}]{61}

que $\mathfrak J^n\subset V$ (ce qu'on exprime, par abus de langage, en disant
que la suite $(\mathfrak J^n)$ tend vers $0$). On dit qu'un anneau linéairement
topologisé $A$ est préadmissible s'il existe dans $A$ un idéal de définition~;
on dit que $A$ est admissible s'il est préadmissible et si en outre il est
séparé et complet.
\end{definition}

Il est clair que si $\mathfrak J$ est un idéal de définition, $\mathfrak L$ un
idéal ouvert de $A$, $\mathfrak J\cap\mathfrak L$ est encore un idéal de
définition~; les idéaux de définition d'un anneau préadmissible $A$ forment
donc un système fondamental de voisinages de $0$.

\begin{lemma}[7.1.3]
\label{0.7.1.3-fr}
Soit $A$ un anneau linéairement topologisé.
\begin{enumerate}
  \item[{\upshape(i)}] Pour que $x\in A$ soit topologiquement nilpotent, il
    faut et il suffit que pour tout idéal ouvert $\mathfrak J$ de $A$, l'image
    canonique de $x$ dans $A/\mathfrak J$ soit nilpotente. L'ensemble
    $\mathfrak T$ des éléments topologiquement nilpotents de $A$ est un idéal.
  \item[{\upshape(ii)}] Supposons en outre que $A$ soit préadmissible, et soit
    $\mathfrak J$ un idéal de définition de $A$. Pour que $x\in A$ soit
    topologiquement nilpotent, il faut et il suffit que son image canonique dans
    $A/\mathfrak J$ soit nilpotente~; l'idéal $\mathfrak T$ est l'image
    réciproque du nilradical de $A/\mathfrak J$ et est donc ouvert.
\end{enumerate}
\end{lemma}

(i) découle immédiatement des définitions. Pour prouver (ii), il suffit de
remarquer que pour tout voisinage $V$ de $0$ dans $A$, il existe $n>0$ tel que
$\mathfrak J^n\subset V$~; si $x\in A$ est tel que $x^m\in\mathfrak J$, on a
$x^{mq}\in V$ pour $q\geq n$, donc $x$ est topologiquement nilpotent.

\begin{proposition}[7.1.4]
\label{0.7.1.4-fr}
Soient $A$ un anneau préadmissible, $\mathfrak J$ un idéal de définition de
$A$.
\begin{enumerate}
  \item[{\upshape(i)}] Pour qu'un idéal $\mathfrak J'$ de $A$ soit contenu
    dans un idéal de définition, il faut et il suffit qu'il existe un entier
    $n>0$ tel que ${\mathfrak J'}^n\subset\mathfrak J$.
  \item[{\upshape(ii)}] Pour qu'un $x\in A$ soit contenu dans un idéal de
    définition, il faut et il suffit qu'il soit topologiquement nilpotent.
\end{enumerate}
\end{proposition}

(i) Si ${\mathfrak J'}^n\subset\mathfrak J$, pour tout voisinage ouvert $V$ de
$0$ dans $A$, il existe $m$ tel que $\mathfrak J^m\subset V$, donc
${\mathfrak J'}^{mn}\subset V$.

(ii) La condition est évidemment nécessaire~; elle est suffisante, car si elle
est remplie, il existe $n$ tel que $x^n\in\mathfrak J$, donc
$\mathfrak J'=\mathfrak J+Ax$ est un idéal de définition, puisqu'il est ouvert,
et que ${\mathfrak J'}^n\subset\mathfrak J$.

\begin{corollary}[7.1.5]
\label{0.7.1.5-fr}
Dans un anneau préadmissible $A$, un idéal premier ouvert contient tous les
idéaux de définition.
\end{corollary}

\begin{corollary}[7.1.6]
\label{0.7.1.6-fr}
Les notations et hypothèses étant celles de (7.1.4), les propriétés suivantes
d'un idéal $\mathfrak J_0$ de $A$ sont équivalentes~:
\begin{enumerate}
  \item[$a)$] $\mathfrak J_0$ est le plus grand idéal de définition de $A$~;
  \item[$b)$] $\mathfrak J_0$ est un idéal de définition maximal~;
  \item[$c)$] $\mathfrak J_0$ est un idéal de définition tel que l'anneau
    $A/\mathfrak J_0$ soit réduit.
\end{enumerate}
Pour qu'il existe un idéal $\mathfrak J_0$ ayant ces propriétés, il faut et il
suffit que le nilradical de $A/\mathfrak J$ soit nilpotent~; $\mathfrak J_0$
est alors égal à l'idéal $\mathfrak T$ des éléments topologiquement nilpotents
de $A$.
\end{corollary}

Il est clair que $a)$ implique $b)$, et $b)$ implique $c)$ en vertu de
(7.1.4, (ii)) et (7.1.3, (ii))~; pour la même raison, $c)$ entraîne $a)$. La
dernière assertion résulte de (7.1.4, (i)) et (7.1.3, (ii)).

Lorsque $\mathfrak T/\mathfrak J$, nilradical de $A/\mathfrak J$, est
nilpotent, on note $A_{\mathrm{red}}$ l'anneau quotient (réduit)
$A/\mathfrak T$.

